% !TEX root = ../main.tex

The locked holdout (19{,}980 searches, 495{,}023 rows) is touched only at this point. The headline number for each predictor is in Table~\ref{tab:hold}.

\begin{table}[t]
\centering
\caption{Holdout NDCG@5 by predictor. The blend selects $w_{\text{xgb}} = 0.5$ on val. The Kaggle public NDCG@5 of the same submission is 0.41706, slightly above the holdout (see Section~\ref{sec:final}).}
\label{tab:hold}
\begin{tabular}{lc}
\toprule
Predictor & Holdout NDCG@5 \\
\midrule
LightGBM seed 42 alone & 0.40789 \\
LightGBM seed 123 alone & 0.40600 \\
LightGBM seed 456 alone & 0.40711 \\
LightGBM 3-seed average & 0.40854 \\
XGBoost \texttt{rank:ndcg} alone & 0.41234 \\
\textbf{Final blend ($w_{\text{xgb}} = 0.5$)} & \textbf{0.41521} \\
\bottomrule
\end{tabular}
\end{table}

The blend on holdout is +0.0029 over XGBoost alone and +0.0067 over the LightGBM 3-seed average, in line with the val numbers. Every predictor sits slightly above its val counterpart, which means the blend choice is not a val-fit artefact. The XGBoost single-model number is above the LightGBM ensemble on both val and holdout, but the blend is consistently above both, confirming that the LightGBM/XGBoost disagreement is the load-bearing source of ensemble lift on this dataset.

\subsection{Segment behaviour}

Search-page size has the largest single effect on conditional NDCG@5 (Fig.~\ref{fig:ndcg-size}, left). Searches with 1 to 10 candidates score a mean NDCG@5 of 0.637, while pages of 31 to 40 candidates drop to 0.347. The bulk of holdout searches sits in the 31 to 40 bucket, which dominates the headline number. Smaller pages saturate NDCG@5 more easily because there are fewer non-positives competing for the top-five slots.

\begin{figure}[t]
\centering
\includegraphics[width=0.85\linewidth]{eval_ndcg_by_size.png}
\caption{Holdout mean NDCG@5 by search-page size. Smaller pages saturate easily because the top-five window is a large share of the candidate set.}
\label{fig:ndcg-size}
\end{figure}

Splitting on \texttt{random\_bool} reveals a 0.086 gap between ranked pages (NDCG@5 = 0.442) and randomly shuffled pages (NDCG@5 = 0.356). The model performs well above a random baseline on shuffled pages, which would score near 0.05, so the lift is real preference learning rather than positional mimicry. The 0.086 gap nonetheless indicates that part of the score on ranked pages comes from features that correlate with Expedia's own ordering (\texttt{prop\_location\_score2} and the historical priors reward the kind of properties Expedia tends to put on top). Booking window has a smaller monotone effect: same-day searches score 0.435 and 90-day-plus searches drop to 0.402, consistent with the intuition that far-in-advance planners click more exploratively and are harder to predict from the current feature set.

\subsection{Worst-case characterisation}

About 40\% of holdout searches score exactly NDCG@5 = 0 even though positives are present in the candidate set. NDCG@5 is binary in spirit on this dataset: either the booked or clicked hotel makes the top five or it does not, and with 95\% of training rows carrying no positive signal there is a hard floor that cannot easily be pushed without per-user or per-session features that the dataset does not expose. Browsing the worst examples, the dominant pattern is "the user clicked or booked an unremarkable property that the model had no signal about", typically a low-popularity \texttt{prop\_id} on a long page. This is consistent with the popularity-tertile NDCG numbers reported in Section~\ref{sec:bias}.

\subsection{Feature importance}

Figure~\ref{fig:imp} reports the top-20 features by gain for both models. LightGBM and XGBoost agree on the rough ordering. The \texttt{prop\_id} categorical dominates both rankings, followed by \texttt{prop\_location\_score2}, the \texttt{prop\_id} historical priors from block C, and the within-search relativisations from block B (\texttt{price\_z\_within\_srch} and \texttt{price\_rank\_within\_srch} both appear in the top 10). The competitor block barely registers on either model, which is consistent with the missingness pattern. The destination priors from block D rank lower than the \texttt{prop\_id} priors, consistent with their cardinality argument: destinations are coarser than properties.

\begin{figure}[t]
\centering
\includegraphics[width=0.95\linewidth]{feature_importance_top20.png}
\caption{Top-20 features by gain for the LightGBM seed-42 anchor (left) and the XGBoost model (right). The two models agree on the dominant role of \texttt{prop\_id} and the location score; within-search relativisations and historical priors fill the next ranks.}
\label{fig:imp}
\end{figure}
